% ================================================================
% ch06_sakhis_appendix.tex — अध्याय ६: प्रामाणिक साखी संग्रह आ परिशिष्ट
% ================================================================

\cleardoublepage
\section*{अध्याय ६: प्रामाणिक साखी संग्रह आ परिशिष्ट}
\addcontentsline{toc}{section}{अध्याय ६: प्रामाणिक साखी संग्रह आ परिशिष्ट}

\subsection*{६.१ प्रामाणिक साखी संग्रह}
\addcontentsline{toc}{subsection}{६.१ प्रामाणिक साखी संग्रह}

{\small\color{labelcolor}
दरिया साहब के साखियाँ उनकर जीवन-अनुभव के निचोड़ बाड़ी। ई छोट-छोट दोहे अपने में पूरा दर्शन समेटले बाड़ें:
}

\vspace{0.8em}

\begin{quote}
\addfontfeatures{Color=3D3428}
\textbf{साखी १:}\\[0.2em]
दरिया लच्छन साध का, क्या गिरही क्या भेख।\\
नि:कपटी निरसंक रहि, बाहर भीतर एक॥
\\[0.8em]
\textbf{साखी २:}\\[0.2em]
दरिया विरही साध का, तन पीला मन सुक।\\
रेन न आवे नींदरी, दिवस न लागे भूख॥
\\[0.8em]
\textbf{साखी ३:}\\[0.2em]
बड के बड लागै नहीं, बड के लागै बीज।\\
दरिया नान्हा होयकर, रामनाम गह चीज॥
\\[0.8em]
\textbf{साखी ४:}\\[0.2em]
प्रेम ज्ञान जब उपजे, चले जगत कंह झारी।\\
कहे दरिया सतगुरु मिले, पारख करे सुधारी॥
\\[0.8em]
\textbf{साखी ५:}\\[0.2em]
माला तोपी भेष नहिं, नहिं सोना श्रृंगार।\\
सदा भाव सत्संग है, जो कोई गहे करार॥
\\[0.8em]
\textbf{साखी ६:}\\[0.2em]
सुमिरन ऐसा कीजिए, दूजा मन नहिं आय।\\
रोम रोम में नाम रमै, सुरति निरति ठहराय॥
\\[0.8em]
\textbf{साखी ७:}\\[0.2em]
भक्ति करे सो सुरमा, तन मन लज्जा खोय।\\
छैल चिकनिया बिसनी, वा से भक्ति न होय॥
\\[0.8em]
\textbf{साखी ८:}\\[0.2em]
दरिया सरल समर्पणा, यही मोक्ष के उपाय।\\
जो कोई इस राह चले, सो जग से उठि जाय॥
\end{quote}

\begin{center}
{\color{bordercolor}\rule{0.4\textwidth}{0.4pt}}
\end{center}

\subsection*{६.२ परिशिष्ट: स्रोत आ प्रामाणिकता}
\addcontentsline{toc}{subsection}{६.२ परिशिष्ट: स्रोत आ प्रामाणिकता}

\textbf{मुख्य स्रोत:}

\begin{itemize}
  \item \textbf{राधास्वामी सत्संग ब्यास {\engfont (RSSB)}} --- {\engfont``Sant Dariya (Bihar Wale)''} ग्रंथ संपादन।
  \item \textbf{हिंदवी / रेख़्ता डिजिटल आर्काइव} --- संत दरिया साहब के प्रामाणिक पद आ दोहा संकलन।
  \item \textbf{कबीर पारख संस्थान} --- धरकंधा (भोजपुर) आ कबीरपंथी मौखिक वाणी परंपरा।
  \item \textbf{सुफिनामा लाइब्रेरी} --- 18वीं सदी के संत-काव्य आ सधुक्कड़ी संग्रह।
\end{itemize}

\textbf{भाषा-शैली टिप्पणी:}

संत दरिया साहब के मूल वाणी में 18वीं शताब्दी के भोजपुरी भाषा के प्राथमिक रूप के साथ-साथ अवधी आ तत्कालीन सधुक्कड़ी बोली के सम्मिश्रण मिलेला। एह संकलन में हर पद के साथ विशुद्ध आधुनिक भोजपुरी भावार्थ आ अंग्रेजी अनुवाद दिहल गइल बा, ताकि भोजपुरी भाषा आ दर्शन के वैश्विक स्तर पर प्रचार-प्रसार हो सके।

\vfill

\begin{center}
{\color{bordercolor}\rule{0.4\textwidth}{0.6pt}}\\[0.6em]
{\small\color{bordercolor}
  भोजपुरी रिवाइव प्रोजेक्ट\\
  {\engfont Bhojpuri Revive Project | 2026}\\
  {\engfont ``Preserving the 18th Century Bhojpuri Literary Heritage''}
}
\end{center}
